% =============================================================================
\subsection{Discussão dos Resultados e Ameaças à Validade}
\label{sec:discussao}

Os resultados obtidos nas duas fases fornecem lições importantes para o
desenvolvimento de algoritmos de Programação Genética para regressão simbólica,
tanto sob a ótica teórica quanto prática.

\subsubsection{Implicações para o Projeto de Algoritmos de PG}

\paragraph{1. O papel da seleção baseada em casos ($\eps$-Lexicase):}
A seleção por torneio e a seleção por roleta agregam o erro em um escalar
médio (RMSE). Essa compressão destrói a informação sobre o comportamento local
dos indivíduos nos diferentes pontos do domínio, gerando atratores precoces e
convergência prematura (como observado pelo pico de 27 clones no torneio).
A $\eps$-lexicase atua como um mecanismo intrínseco de manutenção de nichos,
permitindo que candidatos subótimos na média geral sobrevivam enquanto
forem especialistas em regiões desafiadoras do espaço de entrada. Esse achado
alinha-se diretamente às descobertas recentes da literatura de
PG~\cite{helmuth2015lexicase,lacava2016epsilon}.

\paragraph{2. Desacoplamento estrutural vs. numérico (Linear Scaling):}
Um dos maiores gargalos observados na Fase~1 foi a tentativa inglória da árvore
de sintetizar constantes numéricas através de longas cadeias aritméticas
(ex.: $(0.55/0.32) \times \dots$). O \emph{Linear Scaling}~\cite{keijzer2003improving}
altera fundamentalmente a paisagem: transforma uma busca não-linear conjunta
por topologia e coeficientes em uma busca estrutural pura pela \emph{forma} da
função, delegando os graus de liberdade lineares à álgebra clássica de mínimos
quadrados em tempo $\mathcal{O}(N)$. Esse mecanismo foi o principal
responsável por destravar a evolução no \code{concrete1} ($R^2$ saltando de
$0.44 \to 0.78$) e por estabilizar o \code{synth2}.

\paragraph{3. A falácia da divisão protegida clássica:}
A convenção de Koza~\cite{koza1992genetic} de retornar $1.0$ apenas quando
$|b| < 10^{-6}$ cria falésias e singularidades assintóticas quando
$10^{-6} < |b| \ll 1$. Em problemas com distribuição contínua de teste fora
dos nós de interpolação do treino, essa descontinuidade induz sobreajuste
severo. A substituição pelo \emph{Analytic Quotient} de \citet{ni2013analytic}
demonstrou ser uma das intervenções de melhor relação benefício/custo:
eliminou 100\% das explosões assintóticas sem acrescentar nenhum custo
computacional mensurável.

\paragraph{4. Recombinação macro vs. exploração micro:}
A constatação empírica de que cerca de $89\%$ dos crossovers deterioram o
melhor pai desmistifica a ideia de que o crossover de subárvore atua
predominantemente recombinando ``blocos construtivos'' (\emph{building blocks})
úteis em regressão simbólica, na veradade, a representação em árvore sofre de forte
epistasia posicional. O rebalanceamento dos operadores com mutação mista
(perturbação gaussiana em ERC e mutação pontual) suavizou a paisagem,
transformando a busca em um processo de exploração global intercalado por
gradientes estocásticos locais.

\subsubsection{Ameaças à Validade e Limitações}

\begin{itemize}
  \item \textbf{Custo computacional da $\eps$-Lexicase:}
        Embora a $\eps$-lexicase tenha sido o método mais eficaz, seu custo
        computacional escala com $\mathcal{O}(P \cdot N)$, exigindo a
        avaliação de erros individuais caso a caso. Em problemas com centenas
        de milhares de amostras, variantes como \emph{Down-sampled
        Lexicase}~\cite{hernandez2019downsampled} seriam necessárias.
  \item \textbf{Critério de parada fixo:}
        Todos os experimentos utilizaram um orçamento fixo de 100 gerações.
        Embora tenhamos mostrado que a saturação ocorria por volta da geração
        80 nas execuções de 100 gerações, orçamentos adaptativos com parada
        antecipada baseada em estagnação poderiam economizar tempo de máquina.
  \item \textbf{Divisão fixa Treino/Teste:}
        Em conformidade com as diretrizes do enunciado da disciplina,
        utilizou-se a partição pré-definida de treino e teste, sem validação
        cruzada ($k$-fold). Esse protocolo reflete com precisão cenários reais
        de teste cego, mas os resultados podem apresentar sensibilidade à
        distribuição das partições fornecidas.
\end{itemize}
